\documentclass[a4paper,11pt]{article}

\usepackage{revy}
\usepackage{amssymb}
\usepackage{amsmath}
\usepackage{commath}
\usepackage{amsfonts}

\revyname{MatematikRevyen}
\revyyear{2025}
\version{0.1}
\eta{? minutter}
\status{Færdig}

\title{Flirt med eksamensvagt}
\author{Augusta '20, Hugo '21, Malthe '20}

\begin{document}
\maketitle

\begin{roles}
\role{I}[KE] Instruktør
\role{S}[Kisser] Studerende til eksamen - kæk og selvsikker 
\role{V}[Elinborg] Ven til eksamen - fornuftig
\role{E}[Louie] Eksamensvagt - en sur og træt ældre dame


\end{roles}

\begin{props}
    \prop{1} Stor frixion kuglepen.
\end{props}


\begin{sketch}


\scene To venner står i døren til eksamenslokalet. V har en stor rygsæk på og noter i hånden. S har intet med. To tomme eksamensborde står midt på scenen. Ude i siden af scenen sidder en ældre kvindelig eksamensvagt ved et bord og krydser eksamensnumre/navne af. 

\says{S} FunkAn eksamen bliver mega nem. 

\says{V} Hvad snakker du om? 

\says{S} Jeg har fået over 80 procent i begge afleveringer så hvis jeg bare kan få skrevet et eller andet ned, er den hjemme. 

\says{V} Jaer okay. Det er vel også bare linear algebra med sådan lidt flere dimensioner. Og jeg har taget gode noter undervejs.

\says{S} Har ikke engang taget noter med. Det er bare mig, min hjerne og min kuglepe-

\scene S klapper sig selv på bukselommen, pauser og klapper febrilsk sig selv over det hele.

\says{S} Shit... ej har du en ekstra kuglepen jeg kan låne? 

\says{V} Nej har kun én.. 

\says{S} Kun én?? Er det ikke lidt risky, hvad nu hvis du løber tør for blæk undervejs?

\says{V} Bro, det er dig der har nul med. Du må finde en anden at låne af... 

\scene Begge kigger hen på eksamensvagten som fortsat skriver med en kuglepen.

\says{S} Jeg spørger hende! 

\says{V} Det kan du ikke! Hun hader dig jo på grund af sidste eksamen!

\says{S} Gud ja, det var hende! Hun var så overdramatisk! Det er jo ikke forbudt at spise til eksamen.

\says{V} Jaer okay men det var måske lige i overkanten at få take-away leveret undervejs. Jeg siger det bare, hun kommer aldrig til at hjælpe dig! 

\says{S} Jeg tror du glemmer hvem du taler til. Jeg kan charmere mig ind på ALLE! Jeg kan bruge Rizz's representationsætning. Jeg skal nok få den kuglepen, bare vent. 

\scene V ryster på hovedet. S tager en dyb indånding, går hen mod eksamensvagten, og mumler til sig selv:

\says{S} Uanset hvad det kræver. 

\scene S kommer hen til bordet. Eksamensvagten kigger træt og vredt op. 

\says{E} Dig igen. 

\says{S} Skal du ikke spørge om mit navn?

\says{E} *Dybt suk* Hvad er dit navn?

\says{S} Lige hvad du har lyst til at kalde mig;)

\says{E} Hvad?

\says{S} S-Simon.

\says{E} Nummer?

\says{S} Spørger du om mit nummer allerede, okay. Jeg kan godt lide kvinder der er direkte.

\says{E} Dit eksamensnummer.

\says{S} Det er forresten en virkelig flot strikvest du har på. Farven matcher virkelig... din... frixion kuglepen.

\says{E} Hvad? 

\says{S} Jeg tænkte lidt på om jeg måtte låne den de næste 3 timer? 

\says{E} Min strikvest? 

\says{S} Nej, fjolle. Din kuglepen! 

\says{E} Nej du kan jo se at jeg bruger den.

\scene Der er stille i lang tid mens E og S kigger på hinanden. 

\says{S} Synes du også denne her øjenkontakt er lidt magisk? 

\says{E} Jeg venter stadig på dit eksamensnummer. 

\says{S} Ah klart. 15. (*løber fra bordet med det samme) 

\scene S skynder sig ned på sin plads. Eksamensvagten går op midt på scenen og taler MEGET dårligt engelsk. 

\says{E} Okay everybody. Welcome to eksamen in funktional analysis.  Please write your eksamensnummer and name on the front page. You have 3 hours. Are there any spørgsmål?

\scene S rækker hånden op med det samme. V tager sig til hovedet. E peger på S.

\says{E} Yes?

\says{S} Er du et Hilbert rum? For du er fuldstændigt ... fantastisk.

\says{E} *Perplexed/rødmer* I can not answer... faglige question.

\says{E} Your 3 hours start now.

\scene Eksamen går i gang. V skriver løs mens S sidder og kigger rundt, trommer lidt i bordet fx. Der går lidt tid hvor S og eksamensvagten udveksler blikke - S ser drømmende på eksamensvagten og vinker flirtende.

\scene S rækker hånden op. E går langsomt ned til S med papir i hånden. 

\says{E} Hvad?

\says{S} Måå jeg måske få noget mere papir? 

\scene S tager papir fra E, men fører flirtende hånden ned langs hendes arm og tager fat i papiret. E virker lidt forskrækket.

\says{S} Hvad tid får du fri?

\says{E} Øøhm altså jeg er pensionist.. jeg har altid fri. 

\scene E går væk - begynder at gå frem og tilbage mellem bordene. 
S sidder og tænker lidt. Når E når rundt til S's bord igen sørger S for at falde af stolen lige foran E. ("falde" betyder

\says{S} Hovsa.. jeg tror jeg er... faldet for dig! 

\scene E ryster på hovedet, går tilbage op til sit bord og sætter sig. S ligger stadig på gulvet.  

\says{V} *hvisker højt ned til S* Var det den rizz du snakkede om tidligere? 

\says{S} *Hvisker højt* Det virker, bare vent! 

\scene S kravler tilbage på sin stol og  rækker hånden op igen. E sukker dybt og går tilbage ned til S. 

\says{S} Hvis jeg skal på toilettet, så skal du gå med... ikke?

\scene E ruller med øjnene. Holder en lang pause. Nikker så til sidst. 

\says{E} Jo okay så kom. 

\scene E forlader scenen først. S går efter, danser forbi V som kigger skræmt på ham. 

\scene E tager sin mobil ud af lommen og ringer til Magdalena.

\says{E} Magdalena, du bliver nødt til at hjælpe mig. Sættet er meget sværere, end du sagde det ville blive. 

\scene Kort pause...

\says{E} Hvis ikke du hjælper mig, så skriver jeg speciale hos Ian i stedet.

\scene E kommer tilbage med rodet hår og strikvesten skævt på - retter på tøjet når de kommer ind igen. S kommer tilfreds tilbage med kuglepennen i hånden og vifter den foran V.

\scene E sætter sig på stolen og har tydeligt ondt bagi, rykker lidt rundt på stolen. E tager et stykke papir frem og tørrer noget af sin mund. 

\says{E} Uanset hvad det kræver.

\scene Lys ned.

\end{sketch}

\end{document}


